\chapter{Introduction}
\label{ch:intro}

%-------------------------------------------------------

Filtering is one of the most common operations in image processing. It
is used to remove noise, to smooth an image before further analysis,
and to find edges, which are the starting point of many computer vision
tasks. In this assignment I implemented three classic filters --- the
box filter, the Sobel filter, and the Gaussian filter --- both by hand
and with OpenCV, and applied each of them to the same input image with
$3 \times 3$ and $5 \times 5$ kernels.

%%%%%%%%%%%%%%%%%%%%%%%%%%%%%%%%%%%%%%%%%%%%%%%%%%%%%
\section{Objectives}
\label{sec:objectives}
%%%%%%%%%%%%%%%%%%%%%%%%%%%%%%%%%%%%%%%%%%%%%%%%%%%%%

The assignment requires implementing the following filters, each with a
$3 \times 3$ and a $5 \times 5$ kernel:

\begin{enumerate}
  \item Box filter, written without OpenCV.
  \item Box filter, using OpenCV.
  \item Sobel filter for X-axis edges, written without OpenCV.
  \item Sobel filter for Y-axis edges, written without OpenCV.
  \item Sobel filter combining X-axis and Y-axis edges, written without
        OpenCV.
  \item Sobel filter combining X-axis and Y-axis edges, using OpenCV.
  \item Gaussian filter, using OpenCV.
\end{enumerate}


%%%%%%%%%%%%%%%%%%%%%%%%%%%%%%%%%%%%%%%%%%%%%%%%%%%%%
\section{Background}
\label{sec:background}
%%%%%%%%%%%%%%%%%%%%%%%%%%%%%%%%%%%%%%%%%%%%%%%%%%%%%

In Assignment 1 we worked with point operators, where each output pixel
depends only on the matching input pixel. Filters are
\emph{neighborhood operators}: each output pixel is computed from a
small window of pixels around it \citep{szeliski2022}. The window is
weighted by a small matrix called the \emph{kernel} $h$, and the
operation is written as
\begin{equation}
  g(x,y) = \sum_{i=-r}^{r} \sum_{j=-r}^{r} h(i,j)\, f(x+i,\, y+j)
  \label{eq:correlation}
\end{equation}
where $f$ is the input image, $g$ is the output image, and the kernel
has size $k \times k$ with $r = \lfloor k/2 \rfloor$. The kernel
``slides'' over the image, and at every position the pixels under it are
multiplied by the kernel weights and summed. The filters in this
assignment only differ in the values of $h$.

\subsection{Box Filter}

The box (mean) filter replaces every pixel with the average of its
$k \times k$ neighborhood, so all weights are equal to $1/k^2$
\citep{gonzalez2018}:
\begin{equation}
  h_{3\times3} = \frac{1}{9}
  \begin{bmatrix} 1 & 1 & 1 \\ 1 & 1 & 1 \\ 1 & 1 & 1 \end{bmatrix}
  \label{eq:box}
\end{equation}
Averaging reduces noise, but it also blurs edges and fine detail. A
larger kernel averages more pixels and therefore blurs more.

\subsection{Sobel Filter}

The Sobel filter estimates the image gradient, i.e.\ how fast the
intensity changes in each direction. Its kernel is the product of a
smoothing vector and a derivative vector \citep{szeliski2022}:
\begin{equation}
  K_x = \begin{bmatrix} 1 \\ 2 \\ 1 \end{bmatrix}
        \begin{bmatrix} -1 & 0 & 1 \end{bmatrix}
      = \begin{bmatrix} -1 & 0 & 1 \\ -2 & 0 & 2 \\ -1 & 0 & 1 \end{bmatrix},
  \qquad K_y = K_x^{T}
  \label{eq:sobel3}
\end{equation}
$K_x$ gives the horizontal gradient $G_x$, which is strong on vertical
edges; $K_y$ gives the vertical gradient $G_y$, which is strong on
horizontal edges. The $5 \times 5$ version is built the same way from
the vectors $[1, 4, 6, 4, 1]$ (smoothing) and $[-1, -2, 0, 2, 1]$
(derivative). The two directions are combined into the gradient
magnitude
\begin{equation}
  G = \sqrt{G_x^2 + G_y^2}
  \label{eq:magnitude}
\end{equation}
which responds to edges in any direction.

\subsection{Gaussian Filter}

The Gaussian filter is also a smoothing filter, but instead of equal
weights it uses weights that follow a Gaussian (bell-shaped) curve, so
the center pixel counts the most and far pixels count less
\citep{gonzalez2018}. With $\sigma = 0$, OpenCV derives $\sigma$ from
the kernel size; the resulting 1D weights are $[0.25, 0.5, 0.25]$ for
$3 \times 3$ and $[0.0625, 0.25, 0.375, 0.25, 0.0625]$ for
$5 \times 5$.


%%%%%%%%%%%%%%%%%%%%%%%%%%%%%%%%%%%%%%%%%%%%%%%%%%%%%
\section{Report Organization}
\label{sec:organization}
%%%%%%%%%%%%%%%%%%%%%%%%%%%%%%%%%%%%%%%%%%%%%%%%%%%%%

The remainder of this report is organized as follows. Chapter
\ref{ch:method} describes the development environment and the code for
each of the seven exercises. Chapter \ref{ch:results} presents the
output of every filter on the input image. Chapter \ref{ch:conclusion}
discusses the results and the challenges I faced, and concludes the
report.
